\end{proof}

\begin{proposition}[The special supercharges and the supercharges]
\label{prop:c11-S}
Let $\eta\in\Sigma^*$ and $b\in V$, so that $\rho(b)\eta\in\Sigma$.
\begin{itemize}
\item (i) Graph by graph, the position part along $b$ of the two-body kernel of $(X_a,X_b)$ along $\eta$ is the two-body kernel of $(X_a,X_b)$ along $\rho(b)\eta$; its translation-invariant part is computed with $S^{(1)}(\eta)$ at the spinor $\rho(x-x_O)\eta$, $x_O$ the point of the density, together with the term $-\int(\rho(A|_{W_{\mathbb R}^\perp})\eta,\lambda^*)$.
\item (ii) Under the hypotheses of Proposition \ref{prop:c7-select} and Corollary \ref{cor:c8-homotopy}, on cohomology, the position part along $b$ of $S_1$ along $\eta$ is $Q_1$ along $\rho(b)\eta$ minus the commutator of $\{S^{(1)}(\eta),-\}$ with $P_1$.
\item (iii) On a family on which $P_1\cong0$, in particular, under the hypotheses of Theorem \ref{thm:c11-arity2}, on the families of Definition \ref{def:c11-families} in the translation-covariant scheme at $\theta=\theta_\lambda=\theta_V=-\tfrac13$, the position part of $S_1$ along $\eta$ is $Q_1$ along $\rho(b)\eta$ on cohomology, and $S_1$ along $\eta$ is fixed by $Q_1$ along $\Sigma$ up to the term of Proposition \ref{prop:c7-select} linear in ${\not d}(\rho(x)\eta)=4\eta$. On the scalar bilinears this term is absent (Corollary \ref{cor:c7-S}), and $S_1$ along $\eta$ is, on cohomology, the combination of the functionals of Corollary \ref{cor:c7-S} with the coefficients of $Q_1$ along $\Sigma$ in Corollary \ref{cor:c7-Q}.
\end{itemize}
\end{proposition}

\begin{proof}
(i) By Lemma \ref{lem:c11-lines} (ii), $S^{(1)}(\eta)=\int x^\mu s_\mu+\int s_0$, where $\int s_\mu$ is $S^{(1)}$ along the constant spinor $\rho(e_\mu)\eta\in\Sigma$ and $s_0=-(\rho(A|_{W_{\mathbb R}^\perp})\eta,\lambda^*)$. Insert $x^\mu=x_O^\mu+(x-x_O)^\mu$ and argue as in the proof of Proposition \ref{prop:c11-position} (i): on plane waves the factor $x^\mu$ is $-i\partial/\partial r^\mu$ of the momentum injected at the vertex, and its action on the momentum delta function is $-i$ times the graph with $\int s_\mu$; the same holds for the special edges, the vertices $S^{(1)}(\eta)$ on the cycle and the counterterm vertices, which are the defect insertions of Proposition \ref{prop:c10-arity1}. (ii) Corollary \ref{cor:c8-homotopy} (ii) at $(\eta,P_b)$: by Proposition \ref{prop:c2-special}, $[P_b,\eta]=\rho(b)\eta$, so $[\eta,P_b]=-\rho(b)\eta$, and $S^{(2)}(\eta,P_b)=0$ (Definition \ref{def:c2-S2}). At order $\hbar$ the identity states that the commutator of $S_1$ along $\eta$ with $\{S^{(1)}(P_b),-\}$, plus the commutator of $\{S^{(1)}(\eta),-\}$ with $P_1$, minus $Q_1$ along $\rho(b)\eta$, is a graded commutator with $\Lambda^{(0)}$, which induces $0$ on cohomology; at order $\hbar^0$ the same identity is the statement of (i) for the classical transformation, which fixes the signs. By Proposition \ref{prop:c10-bound1}, $S_1$ along $\eta$ is at most linear in the position, and Lemma \ref{lem:c11-canonical} identifies the first commutator with its position part. (iii) By Proposition \ref{prop:c7-select} the representative of $S_1$ along $\eta$ at a point $x_0$ is a term linear in the value $\rho(x_0)\eta$ plus a term linear in $4\eta$; the first has the position part given by its value at $\rho(b)\eta$, and as $\eta$ and $b$ vary, $\rho(b)\eta$ spans $\Sigma$, so by (ii) it is the representative of $Q_1$ evaluated at $\rho(x_0)\eta$ on cohomology.
\end{proof}

\begin{proposition}[Renormalisation freedom of the odd components and of $K_1$]
\label{prop:c11-freedom}
\begin{itemize}
\item (i) A field redefinition, a closed direction and the outer direction $G_B$ (Proposition \ref{prop:c8-laws} (a), (c)) change each component by minus the bracket with a local functional, a derivation of the letters: they change the one-body kernels and leave the two-body kernels unchanged. $G_B$ changes $Q_1$ along $S_+\otimes W$ by $-a$ and along $S_-\otimes W^*$ by $+a$ times the classical transformation, $S_1$ along $S_+\otimes W^*$ by $+a$ and along $S_-\otimes W$ by $-a$ times it, and leaves $K_1$ unchanged.
\item (ii) The finite counterterms $\hbar F_1^{(0)}$ and $\hbar F_1^{(1)}(x)$ enter the one-body kernels (Remark \ref{rem:c10-scheme1} (a), (b)); the sector coefficients of Definition \ref{def:c11-free} enter the components along $\Delta$ and $P_b$.
\item (iii) A finite part $\psi'$ of the operators changes the component along $x$ by $[\{S^{(1)}(x),-\},\psi']$ (Proposition \ref{prop:c8-laws} (b)). For $\psi'=\tfrac a2\{S^{(1)}(\Delta),-\}$ this is $\tfrac a2\lambda_x\{S^{(1)}(x),-\}$ (Proposition \ref{prop:c9-scheme}): $+a$, $-\tfrac a2$, $+\tfrac a2$ times the classical transformation along $K_b$, $\Sigma$, $\Sigma^*$, a change of the one-body kernels. A finite part $\psi'$ that vanishes on $O$ changes the component of $O$ along $x$ by $-\psi'(\{S^{(1)}(x),O\})$; in this way the finite parts of the operators of the Weyl weight of $\{S^{(1)}(x),O\}$ enter the two-body kernels.
\item (iv) The cutoff shape enters the finite parts of the special-edge cycles (Proposition \ref{prop:c10-shape}, Lemma \ref{lem:c10-special}), the composite pair kernels and the finite vertex \eqref{eq:c10-finitesp}.
\item (v) Under the hypotheses of Theorem \ref{thm:c11-arity2}, in the translation-covariant scheme at $\theta=\theta_\lambda=\theta_V=-\tfrac13$, $P_1\cong0$ on the families of Definition \ref{def:c11-families} (Corollary \ref{cor:c11-D1fam}), and the condition $P_1=0$ of Definition \ref{def:c8-radial} holds there with $\psi'=0$: the components along $K_b$, $\Sigma$ and $\Sigma^*$ are those of the scheme. For equal coefficients $\theta_P\ne\theta$ along $P_b$, the finite part with $a=(\theta-\theta_P)$ units gives $P_1\cong0$ and changes the components along $K_b$, $\Sigma$, $\Sigma^*$ by $a$, $-\tfrac a2$, $\tfrac a2$ times the classical transformation (Corollary \ref{cor:c11-theta} (iv)).
\item (vi) The class $[Q^\xi_1]$ on the cohomology of $Q^\xi_0$, its kernel, image and rank are invariant under (i) and (iii) (Theorem \ref{thm:c8-twist}).
\end{itemize}
\end{proposition}

\begin{proof}
(i) Proposition \ref{prop:c8-laws}; the bracket with a local functional is a derivation in its second argument, so it acts letter by letter. (ii) Remark \ref{rem:c10-scheme1} and Definition \ref{def:c8-components}. (iii) Proposition \ref{prop:c8-laws} (b) and Proposition \ref{prop:c9-scheme}, with $[\{S^{(1)}(x),-\},\psi'](O)=\{S^{(1)}(x),\psi'(O)\}-\psi'(\{S^{(1)}(x),O\})$. (iv) Remark \ref{rem:c10-scheme1} (d). (v) Corollaries \ref{cor:c11-D1fam} and \ref{cor:c11-theta} (iv). (vi) Theorem \ref{thm:c8-twist}.
\end{proof}

\begin{corollary}[$K_1$ in the translation-covariant scheme]
\label{cor:c11-K}
Under the hypotheses of Theorems \ref{thm:c8-adjoint} and \ref{thm:c11-arity2}, on a family of Definition \ref{def:c11-families} at the origin of one length, in the translation-covariant scheme at $\theta=\theta_\lambda=\theta_V=-\tfrac13$, $K_1$ along $K_{\bar b}$ equals $\hbar\big[\{S^{(1)}(K_{\bar b}),-\},\mathrm{Gram}_{\rm tree}^{-1}\mathrm{Gram}_{\rm 1}\big]$, the statement $K=P^\dagger$ of Corollary \ref{cor:c8-gram}. For equal coefficients $\theta_P$ along $P_b$, $K_1$ along $K_{\bar b}$ is $(\theta-\theta_P)$ units times $\{S^{(1)}(K_{\bar b}),-\}$ plus the same commutator, with the radial Gram matrix of Definition \ref{def:c8-inner} in that scheme.
\end{corollary}

\begin{proof}
Corollary \ref{cor:c11-D1fam} gives $P_1\cong0$, resp. $P_1\cong(\theta-\theta_P)$ units times $\{S^{(1)}(P_b),-\}$ with a real coefficient; Corollary \ref{cor:c8-gram} gives the two relations.
\end{proof}

\begin{proposition}[The chiral-primary bilinears along a rank-one supercharge]
\label{prop:c11-plane}
Let $\xi\in S_+\otimes W$ be of rank one, $\xi_\alpha{}^i=u_\alpha z^i$, and let $w_1,w_2\in W_{\mathbb R}^\perp\otimes\mathbb C$ satisfy $\rho(w_1)\xi=\rho(w_2)\xi=0$. Then $(w_1,w_2)=0$, and under the hypotheses of Proposition \ref{prop:c7-select} the component along $\xi$ of $O=\int\omega\,(\iota_{w_1}A,\iota_{w_2}A)$ vanishes at every order at which they hold: its classical part $\{S^{(1)}(\xi),O\}$ and $Q_1(O)$ are $0$, in every scheme in which the counterterms are equivariant for translations, $\mathfrak{so}(V)$, $\mathfrak{so}(W_{\mathbb R}^\perp)$ and $\mathfrak g$.
\end{proposition}

\begin{proof}
$\rho(w_1)\rho(w_2)+\rho(w_2)\rho(w_1)=2(w_1,w_2)$ applied to $\xi\ne0$ gives $(w_1,w_2)=0$. By Corollary \ref{cor:c7-Q}, for $(w_1,w_2)=0$ the component is a multiple of $\int\omega\,\big[(\iota_{w_1}A,(\rho(w_2)\xi,\lambda))+(\iota_{w_2}A,(\rho(w_1)\xi,\lambda))\big]$, which vanishes; the classical component is the same functional (proof of Corollary \ref{cor:c7-Q}).
\end{proof}

\begin{remark}[The pair entries of the holomorphic twist]
